\answer{OR11}
记 $q=1-p$，并令 $u=u_1$。向左每次只走1，因此从 $n$ 到0必须依次经过 $n-1,\ldots,1$。在各下降首达时刻应用强马尔可夫性和平移不变性，得
\[
 u_n=u^n.
\]
这些下降之间允许出现任意多次向右移动，并不要求连续向左。对从1出发的第一步分类，有
\[
 u=p+qu_3=p+qu^3,
\]
即
\[
 (u-1)(qu^2+qu-p)=0.
\]
二次因子的正根为
\[
 r=\frac{\sqrt{1+4p/q}-1}{2}.
\]
若 $p\ge2/3$，则 $r\ge1$，而二次因子的另一根为负，故 $[0,1]$ 内唯一可用的根是1。

若 $p<2/3$，每步的均值为 $2-3p>0$，大数定律表明游走几乎必然趋于正无穷。此时 $u$ 不可能等于1：否则从任意整数到其左邻点的命中概率都是1，反复应用强马尔可夫性便会几乎必然访问所有负整数，与趋于正无穷矛盾。因此必须取 $u=r<1$。综上，
\[
 \boxed{u_n=
 \begin{cases}
 \left(\dfrac{\sqrt{1+4p/(1-p)}-1}{2}\right)^n,
       &0<p<2/3,\\[4pt]
 1,&2/3\le p<1.
 \end{cases}}
\]
首步方程可能有多个概率范围内的根，漂移或最小非负解的性质才决定应选哪一个。
